\section{How the Calendar Was Constructed}

\subsection{A Paschal and weekly core}

The calendar grew from more than one source. The weekly assembly on the Lord's Day and the annual commemoration of the Lord's Passover supplied its fundamental rhythm. The Paschal observance developed into a preparatory fast, the sacred Triduum, an octave, and a fifty-day Easter season ending at Pentecost. Nativity and Epiphany observances, and then their preparatory and festal seasons, supplied a second annual center. These temporal cycles were not assembled from a single late plan; their dates, lengths, and emphases developed over centuries and varied among churches before the Roman forms became stable.

\subsection{Anniversaries, churches, and the fixed cycle}

Local churches also kept the anniversaries of martyrs, commonly at their place of burial. Celebrations of bishops, confessors, virgins, dedications, and Marian mysteries followed. Exchange among local calendars turned some local observances into regional and eventually Roman or universal ones. The sanctoral cycle therefore preserves different historical layers: biblical apostles and martyrs, Roman cults, medieval devotions, founders and missionaries, and saints drawn from an increasingly worldwide Church.

The two cycles interact. Fixed anniversaries can meet a movable Sunday, feria, or solemnity; without rules of precedence the calendar is only an inventory. Roman books progressively supplied ranks and occurrence rules, while local and religious calendars remained additional layers.

\subsection{Codification and accumulated density}

The Roman Breviary of 1568 and Roman Missal of 1570 regularized a received Roman use after the Council of Trent, while allowing sufficiently ancient local uses to remain. Subsequent canonizations, devotional feasts, rank changes, and extensions accumulated. The calendar did not stand still after Trent, and “the Tridentine calendar” is not one immutable list.

Twentieth-century reforms addressed the resulting pressure on Sundays, ferial psalmody, vigils, octaves, and the temporal cycle. Pius~X's 1911 reform rebalanced the Psalter and occurrence; Pius~XII restored the principal Holy Week rites and in 1955 suppressed many vigils and octaves and simplified ranks. John~XXIII's 1960 code reduced the ranks to four classes and supplied the calendar incorporated in the 1962 Missal. Those reforms are direct predecessors, but they are not the governing calendar of this publication.

\subsection{The conciliar criteria}

The Second Vatican Council described the liturgical year as the Church's unfolding of the whole mystery of Christ and required revision so that its traditional character would be preserved while adapted to contemporary conditions. It directed that Sunday and the Proper of Time receive their due precedence, that the faithful's piety toward the saints remain ordered to Christ, and that only saints of genuinely universal importance normally remain in the universal calendar while others could be left to particular churches, nations, or religious families.

These principles explain the architecture of the 1969 reform. The reform was not only deletion or addition. It revised ranks and precedence, restored Sunday and ferial cycles, repositioned some celebrations nearer historically supportable anniversaries, reduced universal duplication, and deliberately relocated many cults to particular calendars. The resulting system retained ancient and inherited material while making explicit choices about universality.

\subsection{Promulgation and continuing amendment}

Paul~VI promulgated the revised norms and General Roman Calendar with \emph{Mysterii Paschalis} on 14~February 1969, effective 1~January 1970. The 1970 Roman Missal embodied the reform; later typical editions consolidated it. The third typical edition of 2002 is the printed base for the fixed inventory here, with its 2008 emended printing and later decrees treated as a continuing amendment chain rather than as a new calendar.

The system remains open to competent amendment. Since 2002, universal entries have been added, one memorial broadened to include a household of three saints, and Mary Magdalene's celebration raised to a feast. In parallel, the United States bishops' approved proper selects national saints, changes several universal dates or ranks within the territory, transfers Epiphany and Corpus Christi, and permits a province-dependent transfer of the Ascension. This continuing life is why the inventory has an exact cutoff.
